\cleardoublepage
\chapternonum{摘要}

随着云计算、大数据与人工智能应用的快速发展，数据中心规模持续扩大。
固态硬盘（SSD）作为关键存储介质，其可靠性与寿命利用效率直接影响业务连续性和运维成本。
本文针对规模化数据中心中SSD故障预警误报成本高、闪存块耐久差异未被充分利用这两大核心问题，以提升运维可靠性和资源利用率为目标，围绕盘级故障预警与块级寿命优化展开系统性研究。

针对现有盘级故障预测方法误报率偏高、难以捕捉近端突变前兆的问题，本文提出一种面向低误报约束的预测方法。
该方法通过混合标注策略与短期时序突变特征，将监督信号与故障前兆的时间尺度对齐，在抑制标签噪声的同时有效放大微弱的故障信号。
模型采用RF、XGB与LGB三种异构基学习器，并引入Attention模块进行样本级自适应加权融合，同时利用Focal Loss缓解类别不平衡，从而提升对稀疏故障样本的判别稳定性。
在阿里云MB1与MB2数据集上的实验结果表明，该方法在维持极低误报率的前提下，分别取得了0.921和0.490的$F_{0.5}$分数，实现了对故障盘的精准、及时检出。

针对传统统一编程/擦除（PE）阈值策略浪费强块寿命的问题，本文设计了一套基于块级错误比特计数（FBC）的差异化寿命提升策略。
该策略将块寿命预测转化为“强/弱块”二分类问题，并据此实施差异化的PE擦写上限管控，在不突破可靠性边界的前提下充分挖掘强块的剩余寿命。
测试结果表明，该策略可带来约17\%的累计可写入量提升。轻量化采集实验进一步验证了该策略在控制器端部署的可行性。

为验证上述算法的工程实用性，本文设计并实现了一套智慧SSD运维原型软件。
该系统打通了从数据采集、模型推理到预警处置的完整业务闭环，并通过功能测试验证了其从理论模型到实际运维流程的转化能力。

本文构建了从“盘级精准预警”到“块级寿命优化”再到“系统闭环落地”的完整技术链路。
研究表明，将低误报约束贯穿于算法设计与策略制定的全流程，能够在保障存储可靠性的同时显著提升资源利用效率。
这一工作为数据中心固态硬盘的可靠性运维与成本优化提供了可行的方法与实践参考。   

\textbf{关键词}：固态硬盘；机器学习；低误报故障预警；寿命提升；

\cleardoublepage

Cloud computing, big data and artificial intelligence are developing rapidly, driving the continuous expansion of datacenter scales. 
Solid State Drives (SSDs) are key storage devices in data centers. 
The stability and effective service life of SSDs greatly affect the normal operation and maintenance cost of data center systems.
This thesis investigates two important problems in SSD daily operation and maintenance. 
The first is SSD failure prediction. 
Early monitoring data of SSDs is usually insufficient and noisy, so this study aims to find faulty SSDs accurately with few false alarms.
The second is SSD service life optimization. Traditional unified service life rules cannot adapt to the different durability of different storage blocks. 
This thesis predicts the service life of each block with almost no false alarms, and designs practical writing rules to improve the total usable writing capacity of SSDs.
This thesis carries out three core research works as follows.

First, this thesis analyzes SSD failure rules based on the public Alibaba Cloud SMART dataset. 
It proposes a new labeling method to reduce data noise in normal working stages. 
It also extracts new time-series features that can capture short-term data changes, and builds a new prediction model. 
The model combines three classic machine learning algorithms (RF, XGB, LGB) and uses an attention mechanism to optimize prediction results. 
All experiments are divided by month in chronological order. 
The model threshold is adjusted and optimized on the verification set, taking disk-level $F_{0.5}$ as the main evaluation standard. 
In three repeated experiments, the model achieves an $F_{0.5}$ of 0.921 and a TPR of 0.738 on the MB1 dataset, and an $F_{0.5}$ of 0.490 and a TPR of 0.187 on the MB2 dataset. 
Meanwhile, the model keeps a low false alarm rate, which proves that it can effectively predict SSD failures with few false alarms.

Second, this thesis uses block-level fault bit count (FBC) data to judge the service life of storage blocks. 
It divides storage blocks into high-quality blocks and low-quality blocks through classification. 
A differentiated writing control strategy is proposed accordingly. 
Strict writing limits are set for low-quality blocks to ensure use safety, while proper higher writing limits are allowed for high-quality blocks on the premise of safety. 
Compared with the traditional unified erasing and writing threshold strategy, this method increases the total usable writing capacity by about 17\% on the held-out test set based on model predictions, and significantly improves storage utilization.

Third, based on the above failure prediction and life optimization methods, this thesis develops an intelligent SSD operation and maintenance software. 
The software realizes full-process management from early fault warning to equipment state monitoring. 
It supports equipment asset management, log recording, fault statistics and early warning processing. 
The practical test proves that the proposed methods can be well applied and deployed in actual engineering scenarios.

In summary, this thesis forms a complete technical solution including failure prediction, service life optimization and system development. 
It improves the accuracy of SSD fault identification and the utilization efficiency of storage resources without increasing false alarms. 
This research provides a reliable and practical reference for daily operation, maintenance and cost control of SSDs in data centers.


\textbf{Keywords}：solid-state drive; machine learning; low false alarm failure prediction; lifetime enhancement; 